\par\addvspace{12pt}\Needspace{9\baselineskip}

\section{函数研究的方法}

% ITEM: calculus-example-045
\solutionitem{calculus-example-045}{例5.45}
设\(f(0)=1\)，且\(x\ge0\)时\(f'(x)\ge f(x)\)。证明\(f(x)\ge e^x\)。再设\(x>0\)时\(xf'(x)-f(x)>0\)，比较\(f(2x)\)与\(2f(x)\)。分别说明为什么选择乘积或商。
\solutionanswer
取\(H=e^{-x}f\)，则\(H'=e^{-x}(f'-f)\ge0\)，故\(H(x)\ge H(0)=1\)，即\(f(x)\ge e^x\)。第二问取\(G=f/x\)，有\(G'=(xf'-f)/x^2>0\)，所以\(f(2x)/(2x)>f(x)/x\)，即\(f(2x)>2f(x)\)。选取乘积或商，是为了让导数中恰好出现题设已控制符号的组合。


% ITEM: calculus-example-046
\solutionitem{calculus-example-046}{例5.46}
比较\(e^\pi\)与\(\pi^e\)的大小。再设\(e<a<b\)，比较\(a^b\)与\(b^a\)。
\solutionanswer
\(H=\ln x/x\)在\((e,+\infty)\)严格递减。由\(\pi>e\)，\(\ln\pi/\pi<1/e\)，故\(e\ln\pi<\pi\)，即\(\pi^e<e^\pi\)。同理\(e<a<b\)给出\(\ln a/a>\ln b/b\)，所以\(a^b>b^a\)。


% ITEM: calculus-exercise-25-01
\solutionitem{calculus-exercise-25-01}{习题5.25（1）}
设\(u,v\in\mathbb R\)，且\(u^3+u=-(v^3+v)\)，证明\(u+v=0\)。
\par
\solutionanswer
取\(H(x)=x^3+x\)，则\(H'=3x^2+1>0\)，且\(H(-v)=-H(v)\)。题设为\(H(u)=H(-v)\)，严格单调性给出\(u=-v\)。


% ITEM: calculus-exercise-25-02
\solutionitem{calculus-exercise-25-02}{习题5.25（2）}
设\(f\)在\([0,1]\)上连续可导，\(f(0)=2\)，且\(f'+2f\le0\)。估计\(f(x)\)，并说明比较方向怎样由所选函数的导数确定。
\par
\solutionanswer
\((e^{2x}f)'=e^{2x}(f'+2f)\le0\)，所以\(e^{2x}f(x)\le f(0)=2\)，即\(f(x)\le2e^{-2x}\)。乘子为正，组合导数的非正性决定上界方向。


% ITEM: calculus-exercise-25-03
\solutionitem{calculus-exercise-25-03}{习题5.25（3）}
设\(f>0\)，且\(f'/f\le1\)于\([a,b]\)。证明\(f(b)/f(a)\le e^{b-a}\)，比较研究\(\ln f\)与\(e^{-x}f\)两条路线。
\solutionanswer
研究\(H=\ln f-x\)，有\(H'\le0\)，故\(\ln f(b)-\ln f(a)\le b-a\)，取指数即得。也可研究\(G=e^{-x}f\)，\(G'=e^{-x}f(f'/f-1)\le0\)，得\(e^{-b}f(b)\le e^{-a}f(a)\)。两种方法分别把相对变化转为差及乘积。


% ITEM: calculus-exercise-25-04
\solutionitem{calculus-exercise-25-04}{习题5.25（4）}
设\(H(0)=H'(0)=H''(0)=0\)，且\(H'''\ge0\)于\([0,1]\)。证明\(H\ge0\)。再将这一推理用于\(e^x-1-x-x^2/2\)，指出每次回推所需的零点信息。
\solutionanswer
由\(H^{(3)}\ge0,H''(0)=0\)得\(H''\ge0\)；再用\(H'(0)=0\)得\(H'\ge0\)；最后用\(H(0)=0\)得\(H\ge0\)。对\(H=e^x-1-x-x^2/2\)，三阶导数\(e^x>0\)，前三个初始值均为零，故结论适用，每次回推都需要对应低一阶函数在起点的值。


% ITEM: calculus-example-047
\solutionitem{calculus-example-047}{例5.47}
设\(f_a(x)=e^x+x^2-ax\)，最小值点为\(c\)。确定\(c>0\)、\(c=0\)、\(c<0\)分别对应的\(a\)的范围，再证明最小值不超过\(1\)，并确定等号条件。比较直接使用\(f_a(0)\)与研究化简后的\(f_a(c)\)两种方法。
\solutionanswer
\(a=e^c+2c\)严格随\(c\)增加，\(c=0\)对应\(a=1\)。所以\(a>1,a=1,a<1\)分别对应\(c>0,c=0,c<0\)。直接比较有\(\min f_a\le f_a(0)=1\)；等号要求零点也是唯一最小值点，故恰为\(a=1\)。也可研究\(M(c)=(1-c)e^c-c^2\)，\(M'=-c(e^c+2)\)，故在零点取唯一最大值\(1\)。前一种比较更短，后一种则给出了最小值随最小值点的变化。


% ITEM: calculus-exercise-26-01
\solutionitem{calculus-exercise-26-01}{习题5.26（1）}
设\(c\)为\(e^x+x-2=0\)的根，证明\(0<c<1\)。利用\(e^c=2-c\)，将\(e^{2c}+2ce^c+c^2\)化为确定的数。
\solutionanswer
\(e^x+x-2\)严格递增，在\(0,1\)的值分别为\(-1,e-1\)，故唯一根在\((0,1)\)。所求式为\((e^c+c)^2=2^2=4\)。


% ITEM: calculus-exercise-26-02
\solutionitem{calculus-exercise-26-02}{习题5.26（2）}
设\(f(x)=x(\ln x-ax)\)有两个极值点\(c_1<c_2\)。由驻点方程化简\(f(c_i)\)，并研究如何利用\(c_i\)的位置估计这两个极值。
\solutionanswer
驻点方程为\(\ln c+1=2ac\)。两个驻点要求\(0<a<1/2\)：\(f'\)在\(1/(2a)\)处取最大值\(\ln(1/(2a))>0\)，两端趋负无穷。且\(e^{-1}<c_1<1<c_2\)，更有\(c_2>1/(2a)\)。代入方程得\(f(c_i)=c_i(\ln c_i-1)/2\)。函数\(J(c)=c(\ln c-1)/2\)在\((0,1)\)减、\((1,+\infty)\)增，故\(-1/2<f(c_1)<-1/e\)，\(f(c_2)>-1/2\)。进一步，\(c_2\)与\(e\)的位置由\(f'(e)=2-2ae\)决定：\(a<1/e\)时极大值为正，\(a=1/e\)时为零，\(1/e<a<1/2\)时为负。


% ITEM: calculus-exercise-26-03
\solutionitem{calculus-exercise-26-03}{习题5.26（3）}
设\(a>0\)，函数\(f(x)=-\ln x-ax^2+x\)有两个极值点\(c_1,c_2\)。先确定\(a\)的范围，再用韦达关系将\(f(c_1)+f(c_2)\)表示为\(a\)的函数。
\solutionanswer
驻点满足\(2ac^2-c+1=0\)。有两个不同正根等价于\(0<a<1/8\)，导数在两根处变号，确为极值点。又\(c_1+c_2=c_1c_2=1/(2a)\)，且\(f(c_i)=-\ln c_i+(c_i+1)/2\)。所以\(f(c_1)+f(c_2)=\ln(2a)+1+1/(4a)\)。


% ITEM: calculus-exercise-26-04
\solutionitem{calculus-exercise-26-04}{习题5.26（4）}
本题\(f_a(x)=e^x+x^2-ax\)，最小值点记为\(c\)。
对正文中的\(f_a\)，若\(a\)只在\([1,e+2]\)内变化，确定最小值点的范围。说明先求\(c\)的范围，再研究\((1-c)e^c-c^2\)，怎样避免同时保留两个独立变量。
\solutionanswer
\(a=e^c+2c\)严格递增，将\([1,e+2]\)对应为\(c\in[0,1]\)。于是最小值为\(M(c)=(1-c)e^c-c^2\)，在此区间从\(1\)严格减至\(-1\)，范围\([-1,1]\)。\(a,c\)由方程相互确定，不应在估计中再把它们当作独立变量。


% ITEM: calculus-example-048
\solutionitem{calculus-example-048}{例5.48}
设\(u<v\)且\(e^u-u=e^v-v\)。令\(d=v-u\)，把\(e^u,e^v\)表示为\(d\)的函数，再研究\(u+v\)。比较差值与比值在本题中的便利程度。
\solutionanswer
由\(e^v-e^u=d\)、\(v=u+d\)，
\[
 e^u=\frac d{e^d-1},\quad e^v=\frac{de^d}{e^d-1},\quad
 u+v=2\ln\frac d{e^{d/2}-e^{-d/2}}<0.
\]
最后一步因为\(e^s-e^{-s}>2s\)对\(s>0\)成立：两边之差在零点为零，导数\(e^s+e^{-s}-2>0\)。该和还随\(d>0\)严格递减，因为其导数为\(2/d-(e^d+1)/(e^d-1)<0\)，这是因为\(J(d)=d(e^d+1)-2(e^d-1)\)满足\(J(0)=J^{(1)}(0)=0\)、\(J^{(2)}(d)=de^d>0\)，逐级回推即得\(J(d)>0\)。两端极限分别为\(0,-\infty\)。本题取差值的好处在于：\(d\)既是两指数值的差，又直接给出它们的比\(e^v/e^u=e^d\)，因此可立即解出两者；知道\(v/u\)则没有这样的关系。


% ITEM: calculus-exercise-27-01
\solutionitem{calculus-exercise-27-01}{习题5.27（1）}
对满足式\textup{(5.24)}的\(u<v\)，令\(t=v/u,\ d=v-u\)，分别用\(t\)、\(d\)表示\(u+v\)与\(uv\)，核对其对应关系。

\par\smallskip\noindent 正文所引关系为：
\begin{equation*}
v-u=\ln\frac vu.
\end{equation*}
\solutionanswer
比值表示给出\(u+v=(t+1)\ln t/(t-1)\)、\(uv=t(\ln t)^2/(t-1)^2\)。差值表示给出\(u+v=d(e^d+1)/(e^d-1)\)、\(uv=d^2e^d/(e^d-1)^2\)。代入\(t=e^d\)两两一致。


% ITEM: calculus-exercise-27-02
\solutionitem{calculus-exercise-27-02}{习题5.27（2）}
对\(0<u<v\)，分别固定\(u\)或令\(t=v/u\)，证明\((v-u)/v<\ln(v/u)<(v-u)/u\)。
\solutionanswer
固定\(u\)，令\(A(v)=\ln(v/u)-(v-u)/v\)，则\(A(u)=0,A'=(v-u)/v^2>0\)；令\(B(v)=(v-u)/u-\ln(v/u)\)，则\(B(u)=0,B'=1/u-1/v>0\)。故两界成立。若令\(t=v/u\)，则只需对\(t>1\)证明\(1-1/t<\ln t<t-1\)，导数判号相同。


% ITEM: calculus-exercise-27-03
\solutionitem{calculus-exercise-27-03}{习题5.27（3）}
设\(0<u<v\)、\(ue^{-u}=ve^{-v}\)。取对数后说明它与式\textup{(5.24)}的关系。

\par\smallskip\noindent 正文所引关系为：
\begin{equation*}
v-u=\ln\frac vu.
\end{equation*}
\solutionanswer
取对数得\(\ln u-u=\ln v-v\)，即\(v-u=\ln(v/u)\)，与所给等值关系完全相同。


% ITEM: calculus-exercise-27-04
\solutionitem{calculus-exercise-27-04}{习题5.27（4）}
设\(u,v>0\)、\(u+v=s\)。用消元、对称性或对数凹性求\(\ln u+\ln v\)的最大值。
\solutionanswer
令\(v=s-u\)，\(0<u<s\)。\(H'=1/u-1/(s-u)\)，在\(s/2\)前正后负，故最大值\(2\ln(s/2)\)，在\(u=v=s/2\)取得。也可直接由对数凹性得到同一结果。


% ITEM: calculus-example-049
\solutionitem{calculus-example-049}{例5.49}
设\(f(x)=xe^{-x}\)，\(u<v\)且\(f(u)=f(v)>0\)。先用反射比较证明\(u+v>2\)，再分别用比值换元与式\textup{(5.25)}证明\(u+v>2\)、\(uv<1\)。比较三条路线所用的数量关系。

\par\smallskip\noindent 正文所引关系为：
\begin{equation*}
\sqrt{ab}<\frac{b-a}{\ln b-\ln a}<\frac{a+b}{2}.
\end{equation*}
\solutionanswer
取对数得\(u-\ln u=v-\ln v\)，故\(0<u<1<v\)。令 \(F(x)=x-\ln x\)。当 \(1<v<2\) 时，写 \(v=1+h\)，\(0<h<1\)。考察反射差
\[
 R(h)=F(1+h)-F(1-h)=2h-\ln\frac{1+h}{1-h}.
\]
因 \(R(0)=0,\ R'(h)=-2h^2/(1-h^2)<0\)，故 \(F(u)=F(v)<F(2-v)\)。左支严格递减，推出 \(u>2-v\)；\(v\ge2\)时结论直接成立。因此\(u+v>2\)。
再令\(t=v/u>1\)，则\(u=\ln t/(t-1)\)、\(v=t\ln t/(t-1)\)。由\(2(t-1)/(t+1)<\ln t<(t-1)/\sqrt t\)，分别得到\(u+v>2\)、\(uv<1\)。比值换元消去了\(u\)，把根和、根积的比较分别化为\(\ln t\)的下界和上界。

若保留两个根，由\(v-u=\ln v-\ln u\)可知，\(u,v\)的对数平均恰为\(1\)。直接用式\textup{(5.25)}，即得\(\sqrt{uv}<1<(u+v)/2\)。反射法比较的是两根离\(1\)的远近，因而只给出根和；后两法用的是同一对数不等式，一种写成比值，另一种直接认出平均值。


% ITEM: calculus-exercise-28-01
\solutionitem{calculus-exercise-28-01}{习题5.28（1）}
若\(u<v\)是\(x-\ln x=a\)的两个根，证明\(a>1\)、\(0<u<1<v\)，并由\(u+v=2a+\ln(uv)\)证明\(u+v<2a\)。
\solutionanswer
\(F=x-\ln x\)在\(1\)处取唯一最小值\(1\)，两根必有\(a>1,0<u<1<v\)。等值关系使对数平均为\(1\)，故\(uv<1\)。将两方程相加得\(u+v=2a+\ln(uv)<2a\)。


% ITEM: calculus-exercise-28-02
\solutionitem{calculus-exercise-28-02}{习题5.28（2）}
在上一题条件下，证明\(u+v>a+1\)。可将目标改写为\(v>e^{1-u}\)，比较\(f(e^{1-u})\)与\(f(u)\)，其中\(f(x)=x-\ln x\)。
\solutionanswer
由\(a=u-\ln u\)，目标等价于\(v>1-\ln u\)，又因\(v=a+\ln v\)，也等价于\(\ln v>1-u\)，即\(v>e^{1-u}\)。令\(r=1-u\in(0,1)\)，则
\(F(e^r)-F(u)=e^r-1+\ln(1-r)<0\)：此差在零点为零，导数\(e^r-1/(1-r)<0\)，因为\(\ln(1-r)<-r\)。\(e^r,v>1\)位于递增支，故\(v>e^r\)，从而\(u+v>a+1\)。


% ITEM: calculus-exercise-28-03
\solutionitem{calculus-exercise-28-03}{习题5.28（3）}
对\(0<a<b\)，用对数平均不等式估计\(\ln b-\ln a\)。比较所得两侧界与中值定理给出的\((b-a)/b\)、\((b-a)/a\)。
\solutionanswer
取倒数并乘\(b-a>0\)，得\(2(b-a)/(a+b)<\ln b-\ln a<(b-a)/\sqrt{ab}\)。左界严格大于\((b-a)/b\)，右界严格小于\((b-a)/a\)，比只用端点导数的中值定理界更紧。


% ITEM: calculus-exercise-28-04
\solutionitem{calculus-exercise-28-04}{习题5.28（4）}
若\(f\)在\(c\)两侧分别严格递增、严格递减，且\(f(c+t)>f(c-t)\)，研究同一高度的两点\(u<c<v\)的中点位于\(c\)的哪一侧。写出所用比较点处于定义域的条件。
\solutionanswer
中点在\(c\)的右侧。若不然，则\(v-c\le c-u\)。把右点\(v\)关于\(c\)反射，得到\(u\le 2c-v<c\)，所以反射点必在左侧定义区间内。取\(t=v-c\)，题设的比较式和左支的单调性给出
\[
 f(v)>f(2c-v)\ge f(u),
\]
与两点同高矛盾。这里要比较的是\(v\)和\(2c-v\)；无需假定\(u\)的反射点也在定义域内。


% ITEM: calculus-exercise-28-05
\solutionitem{calculus-exercise-28-05}{习题5.28（5）}
设\(k>0\)，\(u<v\)满足\(u-k\ln u=v-k\ln v\)。证明\(u+v>2k\)、\(uv<k^2\)，并说明这两个界为什么随\(k\)分别作一次、二次伸缩。
\solutionanswer
令\(u=kr,v=ks\)，原式化为\(r-\ln r=s-\ln s\)。由对数平均关系得\(r+s>2,rs<1\)，还原即\(u+v>2k,uv<k^2\)。和与积分别在长度伸缩下乘\(k,k^2\)。


% ITEM: calculus-exercise-28-06
\solutionitem{calculus-exercise-28-06}{习题5.28（6）}
对\(0<t<1\)，证明\(2t<\ln\frac{1+t}{1-t}<\frac{2t}{\sqrt{1-t^2}}\)，并说明它与反射差\(R(t)\)的关系。
\solutionanswer
对\(a=1-t,b=1+t\)用对数平均不等式，\(\sqrt{1-t^2}<2t/\ln\frac{1+t}{1-t}<1\)，即所求两界。特别\(R(t)=2t-\ln\frac{1+t}{1-t}<0\)，正是反射比较所用的符号。


% ITEM: calculus-exercise-28-07
\solutionitem{calculus-exercise-28-07}{习题5.28（7）}
（2021年新高考I卷第22题第（2）问）设\(a,b\)是两个不相等的正数，且\(b\ln a-a\ln b=a-b\)。证明\(2<1/a+1/b<e\)。
\solutionanswer
令\(u=1/a,v=1/b\)，不妨\(u<v\)。题设化为\(K(u)=K(v)\)，其中\(K(x)=x\ln x-x\)。\(K\)在\(1\)处取最小值，且\(K(0^+)=K(e)=0\)，故\(0<u<1<v<e\)。
先比较\(K(1+t)-K(1-t)\)，其导数为\(\ln(1-t^2)<0\)，所以对\(t=1-u\)，\(K(2-u)<K(u)=K(v)\)。右支递增，推出\(v>2-u\)。
再令\(D(u)=K(e-u)-K(u)\)。\(D(0^+)=0\)、\(D(1)=K(e-1)+1>0\)，且\(D'=-\ln[u(e-u)]\)在\((0,1)\)先正后负，因\(u(e-u)\)严格递增并恰过\(1\)。故\(D(u)>0\)，由右支递增得\(v<e-u\)。综上\(2<u+v<e\)，即原结论。


% ITEM: calculus-example-050
\solutionitem{calculus-example-050}{例5.50 *}
本题的两支根由\(xe^{-x}=a\)、\(0<a<1/e\)确定，且\(0<u(a)<1<v(a)\)。

对上述两支根\(u(a),v(a)\)，研究当\(a\to0^+\)及\(a\to(1/e)^-\)时它们各自的极限。再研究\(v(a)/u(a)\)的单调性，比较直接求商的导数与先求对数导数的方法。
\solutionanswer
这里\(u,v\)满足\(xe^{-x}=a\)，\(0<a<1/e\)，且\(0<u<1<v\)。由两个单调分支的反函数及其端点值，当\(a\to0^+\)时\(u\to0,v\to+\infty\)；当\(a\to(1/e)^-\)时两者均趋\(1\)。由\(x'=x/[a(1-x)]\)，
\[
 \left(\ln\frac vu\right)'=\frac1a\left(\frac1{1-v}-\frac1{1-u}\right)<0.
\]
故\(v/u\)严格递减，从\(+\infty\)趋向\(1\)。直接求商会多带正因子\(v/u\)，对数导数先约掉这个因子，符号更清楚。


% ITEM: calculus-exercise-29-01
\solutionitem{calculus-exercise-29-01}{习题5.29（1） *}
设\(u(a)<v(a)\)为\(x-\ln x=a\)的两个根。求两支的导数，判断根的间距\(v-u\)怎样随\(a\)变化。
\solutionanswer
\(u'=u/(u-1)<0\)，\(v'=v/(v-1)>0\)。所以\((v-u)'>0\)，两根间距严格随\(a>1\)增加。


% ITEM: calculus-exercise-29-02
\solutionitem{calculus-exercise-29-02}{习题5.29（2） *}
证明上题中的\(u+v\)随\(a>1\)增大而增大，并研究\(a\to1^+\)时的极限。与正文中的参数\(a\)比较：两个问题中的参数是否表示同一个量？
\solutionanswer
\((u+v)'=(2uv-u-v)/[(u-1)(v-1)]>0\)，因为\(uv<1,(u+v)>2\)，分子、分母均负。当\(a\to1^+\)，两个反函数分支均趋\(1\)，故和趋\(2\)。本题参数是\(x-\ln x\)的水平高度；正文\(xe^{-x}=a\)的参数则是其指数变换，二者满足\(a_{\text{正文}}=e^{-a_{\text{本题}}}\)，方向相反。


% ITEM: calculus-exercise-29-03
\solutionitem{calculus-exercise-29-03}{习题5.29（3） *}
对\(a=e^c+2c\)确定的\(c=c(a)\)，证明\(c'(a)=1/(e^c+2)\)。再对最小值\(m(a)=e^{c(a)}+c(a)^2-ac(a)\)求导，观察驻点方程怎样使含\(c'\)的项消去。
\solutionanswer
反函数求导得\(c'=1/(e^c+2)\)。直接求导\(m'=(e^c+2c-a)c'-c=-c\)，含\(c'\)的系数恰是驻点方程，因而为零。特别最小值在\(a<1\)时增、\(a>1\)时减。


% ITEM: calculus-example-051
\solutionitem{calculus-example-051}{例5.51 *}
设\(u<v\)为\(x-\ln x=a\)的两个根。取\(q(x)=(x+x^{-1})/2\)，证明它在两根所在的两段上分别位于\(f\)的上、下方，再由\(q(x)=a\)的根证明\(u+v<2a\)。与利用\(uv<1\)的证明比较。
\solutionanswer
令\(D=q-f=\ln x+(x^{-1}-x)/2\)，则\(D(1)=0\)、\(D'=-(x-1)^2/(2x^2)\)。故左支\(q>f\)，右支\(q<f\)。\(q=a\)的两根为\(r_\pm=a\pm\sqrt{a^2-1}\)。\(q\)在\((0,1)\)递减、\((1,+\infty)\)递增；由\(q(u)>a\)、\(q(v)<a\)分别得到\(u<r_-\)、\(v<r_+\)，所以\(u+v<2a\)。另一法由两根方程相加得\(u+v=2a+\ln(uv)\)，再用\(uv<1\)即可；比较曲线法则用一个可显式解的方程约束两个根。


% ITEM: calculus-exercise-30-01
\solutionitem{calculus-exercise-30-01}{习题5.30（1） *}
设\(Q(x)=(1+Ax)/(1+Bx)\)。使它在零点匹配\(e^x\)直到二阶变化，求\(A,B\)。再在\((-2,2)\)内研究\(e^x\)与\(Q(x)\)的大小关系，可以先比较\(x\)与\(\ln Q(x)\)。
\solutionanswer
展开\(Q=1+(A-B)x+(B^2-AB)x^2+o(x^2)\)，匹配得\(A=1/2,B=-1/2\)。在\((-2,2)\)，\(Q>0\)，且\((\ln Q-x)'=x^2/(4-x^2)\ge0\)。零点处差为零，故\(x>0\)时\(Q>e^x\)，\(x<0\)时\(Q<e^x\)，零点相等。


% ITEM: calculus-exercise-30-02
\solutionitem{calculus-exercise-30-02}{习题5.30（2） *}
在\(x=1\)附近比较\(\ln x\)与\(2(x-1)/(x+1)\)的Taylor多项式。两者从哪一阶开始不同？这能否直接代替正文中全区间的误差符号证明？
\solutionanswer
令\(h=x-1\)，对数为\(h-h^2/2+h^3/3+o(h^3)\)，有理式为\(h-h^2/2+h^3/4+o(h^3)\)。从三阶开始不同。局部三次主部只能判断充分靠近\(1\)时的符号；全区间结论仍需正文的导数判号。


% ITEM: calculus-exercise-30-03
\solutionitem{calculus-exercise-30-03}{习题5.30（3） *}
设\(f,g\)在\(x>c\)上均严格递增，且\(g(x)>f(x)\)。若\(r,s>c\)且\(f(r)=g(s)\)，判断\(r,s\)的大小。再写出严格递减时的结论，用以核对正文两支比较的方向。
\solutionanswer
严格递增时\(f(r)=g(s)>f(s)\)，故\(r>s\)。若两者均严格递减，同一比较推出\(r<s\)。函数的上下关系相同，分支的增减方向改变根的位置方向。


% ITEM: calculus-exercise-30-04
\solutionitem{calculus-exercise-30-04}{习题5.30（4） *}
在\(x=0\)处，用\((1+Ax)/(1+Bx)\)匹配\(\sqrt{1+x}\)直到二阶变化。研究所得有理函数与\(\sqrt{1+x}\)在\(x\ge0\)时的大小关系，并估计\(\sqrt{1.1}\)。
\solutionanswer
匹配\(1+x/2-x^2/8\)得\(A=3/4,B=1/4\)，即\(Q=(4+3x)/(4+x)\)。对\(x>0\)，
\[
 Q^2-(1+x)=-\frac{x^3}{16(1+x/4)^2}<0,
\]
所以\(Q<\sqrt{1+x}\)。差为\(x^3/[16(1+x/4)^2(\sqrt{1+x}+Q)]\le x^3/32\)。因此\(\sqrt{1.1}\approx43/41\approx1.04878049\)，近似偏小，误差小于\(1/32000\)。


% ITEM: calculus-exercise-30-05
\solutionitem{calculus-exercise-30-05}{习题5.30（5） *}
（2018年全国I卷理科第21题第（2）问）设\(f(x)=1/x-x+a\ln x\)有两个极值点\(x_1,x_2\)。证明\(\bigl(f(x_1)-f(x_2)\bigr)/(x_1-x_2)<a-2\)。可先利用驻点方程，再选择对数的有理界。
\solutionanswer
驻点方程为\(x^2-ax+1=0\)，故\(a>2\)。不妨\(x_1=u<1<v=x_2\)，\(uv=1\)。在驻点处\(f(x)=a-2x+a\ln x\)，于是
\[
 \frac{f(u)-f(v)}{u-v}=-2+a\frac{\ln v-\ln u}{v-u}.
\]
对\(v>1\)，\((v-v^{-1})/2-\ln v\)的导数为\((v-1)^2/(2v^2)>0\)，故\(2\ln v<v-v^{-1}=v-u\)。代入上式即严格小于\(a-2\)。


% ITEM: calculus-exercise-30-06
\solutionitem{calculus-exercise-30-06}{习题5.30（6） *}
（2021年新高考I卷第22题再研究）回看习题5.28第（7）题。令\(u=1/a,v=1/b\)，不妨设\(u<v\)，再令\(t=v/u\)。将\(u+v\)表示为\(t\)的函数，用本小节的\(\ln t>2(t-1)/(t+1)\)研究其单调性与端点极限，重新证明原结论。
\solutionanswer
变换后\(u\ln u-u=v\ln v-v\)。令\(t=v/u>1\)，可得\(u=\exp(1-t\ln t/(t-1))\)，故
\[
 S(t)=u+v=(1+t)\exp\left(1-\frac{t\ln t}{t-1}\right).
\]
\((\ln S)'=[\ln t-2(t-1)/(t+1)]/(t-1)^2>0\)，所以\(S\)严格递增。由\(\ln t/(t-1)\to1\)得\(S(1^+)=2\)；又\(S=e(1+1/t)\exp[-\ln t/(t-1)]\to e\)，故\(2<S<e\)。与反射法相比，此法把两根和直接变成一个变量的函数。
\par
